\scanpage{316}
\mspage{266}

fonctorielle d'une flèche quelconque (ici les $f_i$) en $pi$, avec $i$ imm ouverte de Dwyer, $p\in\unsure{W_1|W,}$ (p. 229) on est ramené au cas où les $f_i$ sont des imm. ouvertes de Dwyer. \unsure{So} Mais pour \deleted{\unsure{toute|de t}} flèche

% slanted
\margin{C'est un peu vif, car il n'est pas clair que $X'_1\amalg_{X_0}X'_2\longrightarrow X_1\amalg_{X_0}X_2$ soit $\in W$, sans utiliser déjà \deleted{\unsure{la propriété}} \inserted{\unsure{disons}} \unsure{l'un} des trois \unsure{ingrédients} \unsure{$\bigcirc$} de $W(7)$ (savoir $f_0,f_1,f_2\in W\Longrightarrow f_3\in W$). Il semble donc que notre argument ici, \unsure{n'aurait lieu} que moyennant la mise \unsure{hors} \deleted{\unsure{déjà}} \unsure{faite} d'une partie des \unsure{cas admis} \unsure{\ldots} Mais il y a le ``\uline{lemme de Brown} '', cf \uline{XVI} \unsure{\S 3|\S 1}, \unsure{cet argument|ce qui suit} p. \unsure{2|1} \unsure{ff}.}

\inserted{$g:S\mathbin{\unsure{-|\to}}T$} qui est dans $\mathrm{Cof}_W$, soit $Z(g)$ \unsure{son|sa} \unsure{mapping-cone}

\diagram{\(\begin{tikzcd}[ampersand replacement=\&,column sep=small,row sep=small]
S\arrow[rr,"g"]\arrow[dr]\&{}\&T\arrow[dr,"{\alpha(g)}"]\&{}\\
{}\&e\arrow[rr,"{\overline g}"]\&{}\&Z(g)
\end{tikzcd},\)}

\[
\alpha(g):T=\mathrm{but}(g)\longrightarrow Z(g)
\]

la flèche canonique. Celle-ci est fonctorielle en $g$, de sorte que le diagramme cubique donné se prolonge en un double cube

\diagram{\(\begin{tikzcd}[ampersand replacement=\&,row sep=small,column sep=small]
{}\&X_2\arrow[rr]\arrow[dd,dashed,"{f_2}" near end]\&{}\&X_3\arrow[dd,"{f_3}"]\\
X_0\arrow[ur]\arrow[rr]\arrow[dd,"{f_0}"']\&{}\&X_{\unsure{2|1}}\arrow[ur]\arrow[dd,"{f_1}"']\&{}\\
{}\&Y_2\arrow[rr,dashed]\arrow[dd,dashed,"{\alpha_2}" near start]\&{}\&Y_3\arrow[dd,"{\alpha_3}"]\\
Y_0\arrow[ur]\arrow[rr]\arrow[dd,"{\alpha_0}"']\&{}\&Y_1\arrow[ur]\arrow[dd,"{\alpha_1}"']\&{}\\
{}\&Z_2\arrow[rr,dashed]\&{}\&Z_3\\
Z_0\arrow[ur]\arrow[rr,line width=0.9pt]\&{}\&Z_1\arrow[ur]\&{}
\end{tikzcd}\)}

avec \deleted{$\unsure{Z}$} $Z_i=Z(f_i)$, $\alpha_{\unsure{i}}=\alpha(f_i)$ pour $i\in\{0,1,2,3\}$.

Notons que $W(7\ \mathrm{bis})$ implique que si $g\in\mathrm{Cof}_W$,
% archive sheet number: 316
